\documentclass[10pt]{article}
\usepackage[T1]{fontenc}
\usepackage[utf8]{inputenc}
\usepackage[french]{babel}
\usepackage{amsmath,amssymb,mathrsfs,mathpazo}
\usepackage{geometry}
\geometry{a4paper,margin=2cm}
\usepackage{graphicx,array,tabularx,tabularray,colortbl,multirow,multicol,booktabs,enumitem,pifont,ulem,fancybox,diagbox,cancel}
\usepackage{pstricks,pstricks-add,pst-plot,pst-tree,pst-node,pst-text,pst-func,pst-math,pst-bspline,pst-3dplot,pst-eucl,pst-solides3d}
\let\pscalculate\undefined
\usepackage{pst-bezier}
\usepackage{tikz,pgfplots,tkz-tab}
\usetikzlibrary{arrows,arrows.meta,patterns,calc,positioning,trees,decorations.markings,shapes.geometric,shapes.misc,shapes.arrows,intersections,decorations.pathmorphing,decorations.pathreplacing,backgrounds,shadows,babel}
\pgfplotsset{compat=1.18}
\usepgfplotslibrary{fillbetween}
\usepackage[np]{numprint}
\usepackage{listings,makecell,pas-tableur,twoopt}
\let\barre\undefined
\usepackage[pstricks]{tablvar}
\definecolor{darkgreen}{RGB}{0,100,0}
\providecommand{\red}{\color{red}}
\providecommand{\blue}{\color{blue}}
\providecommand{\green}{\color{green}}
\providecommand{\black}{\color{black}}
\providecommand{\white}{\color{white}}
\providecommand{\gray}{\color{gray}}
\providecommand{\cyan}{\color{cyan}}
\providecommand{\magenta}{\color{magenta}}
\providecommand{\yellow}{\color{yellow}}
\providecommand{\violet}{\color{violet}}
\providecommand{\orange}{\color{orange}}
\providecommand{\euro}{\texteuro}
\providecommand{\EUR}{\texteuro}
\providecommand{\ieme}{\textsuperscript{e}}
\providecommand{\degres}{\ensuremath{^\circ}}
\providecommand{\degree}{\ensuremath{^\circ}}
\providecommand{\vv}[1]{\overrightarrow{#1}}
\providecommand{\Coord}[1]{#1}
\providecommand{\decosix}{}
\providecommand{\starredbullet}{\ensuremath{\bullet}}
\pagestyle{empty}
\setlength{\parindent}{0pt}
\usepackage{adjustbox}\DeclareUnicodeCharacter{2236}{:}\DeclareUnicodeCharacter{2190}{\ensuremath{\leftarrow}}
\providecommand{\R}{\mathbb{R}}
\providecommand{\N}{\mathbb{N}}
\providecommand{\Z}{\mathbb{Z}}
\providecommand{\Q}{\mathbb{Q}}
\providecommand{\D}{\mathbb{D}}
\providecommand{\C}{\mathbb{C}}
\providecommand{\e}{{\mathrm{e}}}
\providecommand{\vect}[1]{\overrightarrow{#1}}
\providecommand{\vectt}[1]{\overrightarrow{\mathrm{#1}}}
\providecommand{\vvt}[1]{\overrightarrow{\mathrm{#1}}}
\providecommand{\np}[2][]{{#2~#1}}
\providecommand{\vv}[1]{{\overrightarrow{#1}}}
\renewcommand{\d}{{\mathrm{d}}}
\providecommand{\degres}{{^\circ}}
\providecommand{\degre}{{^\circ}}
\providecommand{\blue}{\color{blue}}
\providecommand{\red}{\color{red}}
\providecommand{\enskip}{\quad}
\providecommand{\llbracket}{\lbrack\!\lbrack}
\providecommand{\rrbracket}{\rbrack\!\rbrack}
\providecommand{\euro}{{\text{€}}}
\providecommand{\no}{{\text{n°}}}
\providecommand{\interval}[2]{{[#1;#2]}}
\providecommand{\intervFO}[2]{{[#1;#2[}}
\providecommand{\dotfill}{\ldots}

\providecommand{\ds}{\displaystyle}
\providecommand{\textding}[1]{#1}
\providecommand{\Oij}{\(\mathrm{O};\vec{\imath},\vec{\jmath}\)}
\providecommand{\Oijk}{\(\mathrm{O};\vec{\imath},\vec{\jmath},\vec{k}\)}
\providecommand{\Ouv}{\(\mathrm{O};\vec{u},\vec{v}\)}
\providecommand{\intervFF}[2]{\ensuremath{[#1;#2]}}
\renewcommand{\euro}{\ifmmode\text{\texteuro}\else\texteuro\fi}
\DeclareUnicodeCharacter{00B0}{\ifmmode^{\circ}\else\textdegree\fi}
\DeclareUnicodeCharacter{20AC}{\ifmmode\text{\texteuro}\else\texteuro\fi}
% Énoncés : Éducation nationale. Transcription et corrigés : APMEP, Denis Vergès et contributeurs. Découpage et classement : CoopMaths, Sébastien Lozano, Éric Elter et contributeurs. Adaptation de la présentation : Maths Explorer.
% CC BY-SA : https://coopmaths.fr/www/about/
% https://coopmaths.fr/alea/static/bac/2023/tex/bac_2023_03_sujet1_polynesie_2.tex
\begin{document}
\noindent\textit{APMEP / CoopMaths. CC BY-SA. Présentation adaptée par Maths Explorer.}\par\medskip
\medskip

L'espace est muni d'un repère orthonormée \Oijk.

On considère:
\begin{itemize}
\item[$\bullet~~$]$d_1$ la droite passant par le point H(2~;~3~;~0) et de vecteur directeur $\vect{u}\begin{pmatrix}1\\-1\\1\end{pmatrix}$ ;
\item[$\bullet~~$]$d_2$  la droite de représentation paramétrique :
\[\left\{\begin{array}{l c l}
x &=& 2k - 3\\
y &=&k\\
z &=&5
\end{array}\right. \quad \text{où }\:k \:\text{décrit }\: \R.\]
\end{itemize}

\medskip

Le but de cet exercice est de déterminer une représentation paramétrique d'une droite $\Delta$ qui soit perpendiculaire aux droites $d_1$ et $d_2$.

\medskip

\begin{enumerate}
\item
	\begin{enumerate}
		\item Déterminer un vecteur directeur $\vect{v}$ de la droite $d_2$. 
		\item Démontrer que les droites $d_1$ et $d_2$ ne sont pas parallèles. 
		\item Démontrer que les droites $d_1$ et $d_2$  ne sont pas sécantes. 
		\item Quelle est la position relative des droites $d_1$ et $d_2$ ?
	\end{enumerate}
\item 
	\begin{enumerate}
		\item Vérifier que le vecteur $\vect{w}\begin{pmatrix}-1\\2\\3\end{pmatrix}$ est orthogonal à $\vect{u}$ et à $\vect{v}$.
		\item On considère le plan $P$ passant par le point H et dirigé par les vecteurs $\vect{u}$ et $\vect{w}$.
		
On admet qu'une équation cartésienne de ce plan est :
\[5x + 4y - z - 22 = 0.\]

Démontrer que l'intersection du plan $P$ et de la droite $d_2$ est le point M(3~;~3~;~5).
	\end{enumerate}
\item Soit $\Delta$ la droite de vecteur directeur $\vect{w}$ passant par le point M. 

Une représentation paramétrique de $\Delta$ est donc donnée par:

\[\left\{\begin{array}{l c l}
x &=&- r + 3 \\
y &= &2r + 3\\
z &=& 3r + 5
\end{array}\right. \quad \text{où }\:r\: \text{décrit }\R.\]

	\begin{enumerate}
		\item Justifier que les droites $\Delta$ et $d_1$ sont perpendiculaires
en un point L dont on déterminera les coordonnées.
		\item Expliquer pourquoi la droite $\Delta$ est solution du problème posé.
	\end{enumerate}
\end{enumerate}

\bigskip

\end{document}
